%---------------------------------------------------------------------
\chapter{Implementation}
\label{ch:impl}
%---------------------------------------------------------------------

\section{Technology Stack}
The backend is a Python package, \code{recycleguard}, built with hatchling and managed with \code{uv}. Its
public surfaces are a Typer command-line interface and a FastAPI application. Persistent state is held in
SQLite databases and in a content-addressed blob store on disk. The agent workflow uses LangGraph
\cite{langgraph} with its SQLite checkpointer. PDF text and layout are extracted with pdfplumber, and all
boundary objects are Pydantic models, from which JSON schemas are exported and checked in. The model is
served by MLX-LM \cite{mlxlm} on loopback and is called over HTTP. The browser interface is a React and
TypeScript application built with Vite. Table~\ref{tab:env} lists the versions.

Local inference dependencies are isolated in a separate project directory, so that the base package, its
API and its tests stay portable and do not require MLX. The model server is started from that directory with
a pinned set of flags, and its session record is hashed into the configuration pin.

\section{Data Preparation and Indexing}
\label{sec:dataprep}
\subsection{Synthetic Dossier Generator}
Dossiers are produced by a generator that builds a typed world model of firms, registrations, obligations,
certificates, transfers and payments with exact decimal quantities, and then renders the visible files.
Faults are injected by named recipes, one per fault family. The generator writes the visible dossier into
Git for the development split and writes the hidden reference records into an evaluator root that is not
shared with the investigator. Construction checks fail the build when splits share an entity, template,
combination, timeline group or world identifier, when a visible name carries a family or seed word, or when a
development and a held-out dossier are too similar (Section~\ref{sec:dataset}).

\subsection{Corpus Indexing}
The regulatory sources are acquired as originals, hashed and retained. Selected provisions are extracted with
their physical page locators and attached to a source version and a reviewed temporal assertion. An FTS5 index is
built over the provision text and is itself stored as an immutable, hashed artifact. A frozen-corpus audit
verifies, before a run, that the pins of the corpus, profile and index match their stored bytes and that the
audit ledger links hold. This audit reports a verified-integrity status only after those checks, and it states
separately that integrity checking does not resolve legal applicability.

\section{Core Agent Loop}
\label{sec:loop}
The workflow is a LangGraph state graph with one node per stage and a conditional edge from the start that
routes to the node named by the current status. Every node ends the graph, so each call to the runner
executes exactly one turn.

\begin{lstlisting}[language=Python,caption={Graph definition (abridged from \code{workflow/runner.py})},label=lst:graph]
builder = StateGraph(GraphInput)
for name in ("preflight", "investigate", "draft", "verify", "sync"):
    builder.add_node(name, self._node)
    builder.add_edge(name, END)
builder.add_conditional_edges(
    START, lambda state: state["route"],
    {name: name for name in ("preflight", "investigate", "draft", "verify", "sync")},
)
self.graph = builder.compile(checkpointer=self.checkpointer)
\end{lstlisting}

Within the \emph{investigate} node the loop is the following. It is written here as pseudocode that follows
the implementation; the code adds the audit records and the error handling.

\begin{algorithm}[H]
\caption{One investigation turn}
\label{alg:turn}
\begin{algorithmic}[1]
\State reserve the token budget for one model call; stop if any limit would be exceeded
\State \textbf{PLAN}: send pending findings, allowed tools and prior observations to the model
\State validate the reply against the plan schema; \textbf{retry once} on invalid output
\If{the reply is a stop}
    \State record the stop reason and move to drafting
\Else
    \State check that the operation is in the whitelist and that its fingerprint is new
    \State store the plan record and the action; execute the tool; store the observation
    \State \textbf{ASSESS}: send plan, action and observation to the model; validate the reply
    \If{the disposition is resolve} \State mark the finding resolved (requires verdict supports)
    \ElsIf{the disposition is abstain} \State mark the finding abstained with its reason
    \Else \State keep the finding pending for the next turn
    \EndIf
    \State pass any finding proposal through the finding gate; retain rejections
\EndIf
\State commit the turn manifest with a compare-and-swap version
\end{algorithmic}
\end{algorithm}

Two repair mechanisms were added after the first real-model runs. If an adaptive action names an observation
identifier that is not in the case, the action is never dispatched or corrected. It is recorded as a rejected
action, and the model is asked once more with the reason, at the cost of one call from the same budget. A second
refusal fails the run. A model reply that fails validation is retried once with a note that the previous reply
was invalid. Both repairs are counted in the evaluation tables so that a reader can see how often they were
needed.

\section{Tools and Function Definitions}
\label{sec:tools}
\begin{table}[H]
\centering
\caption{The six agent tools}
\label{tab:tools}
\small
\begin{tabularx}{\linewidth}{@{}p{4.6cm} X X@{}}
\toprule
\textbf{Tool} & \textbf{Arguments} & \textbf{Result and purpose} \\ \midrule
\code{inspect\_evidence} & Observation identifiers, or empty for a case summary & Parsed fields and exact locators of observations; marks the content as untrusted \\
\code{lookup\_locator} & Source alias or version and a legal path & Verified provision text and metadata at an exact citation \\
\code{trace\_certificate} & Certificate identifiers & Lifecycle (issued, transferred, cancelled, expired); detects reuse and capability mismatch \\
\code{compute\_reconciliation} & Optional bucket identifiers & Required, supported, signed gap and pending quantity per bucket, in exact decimals \\
\code{compare\_return} & Optional bucket identifiers & Compares portal-reported claims with ledger-supported balances \\
\code{search\_rules} & Query and top-$k$ & BM25 passages with authority, temporal and relevance grades \\
\bottomrule
\end{tabularx}
\end{table}

Every tool is read-only. The runner holds the whitelist, so a model cannot add a tool by naming it. The
arguments are typed Pydantic models, which means an out-of-range identifier or a malformed argument is a
validation failure before dispatch and not a runtime surprise.

\section{Prompt Design}
\label{sec:prompt}
The system prompt tells the model that the user message is data and never instructions, that it must return a
single JSON object matching the supplied schema, that it must choose operations only from the allowed tool
names, and that it must abstain when the evidence does not support a conclusion. Each phase adds its own
rules, and the JSON schema of the expected reply is appended to the prompt. The prompt is assembled in the
ChatML format with an empty reasoning block pre-filled, so that the model replies directly with the JSON object
and does not spend tokens on a visible chain of thought. Dossier text reaches the model as escaped JSON, with
angle brackets written as Unicode escapes.

The prompt went through nine versions. The versions after the first real-model run were changed one at a
time, only in response to failures observed on development cases, and each version's results were kept
(Section~\ref{sec:promptdev}). The full text of the final version is in Appendix~\ref{app:prompts}.

\section{Guardrails and Loop Bounds}
\begin{itemize}[itemsep=2pt]
\item \textbf{Tool whitelist.} Six read-only operations; anything else ends the run in an expert-review state.
\item \textbf{Budgets.} For real-model runs: 12 model calls, 600\,s wall time, and a token reservation per call
      and in total, checked before dispatch. The budget is fixed for evaluation and cannot be overridden per run.
      Scripted tests may override it.
\item \textbf{Per-finding limits.} At most eight tool calls, two reformulations, one retry per operation and one
      report repair for a finding or report.
\item \textbf{Cycle detection.} A fingerprint of the finding, tool and arguments stops an exact repeat.
\item \textbf{Finding gate.} Citations must be evidence of this run; the text may contain no digits, prohibited
      claims or authority wording; and a clean finding requires a reconciliation from this run.
\item \textbf{Prohibited outputs.} Legal determination of compliance, allegation of fraud or intent, and penalty
      amounts are rejected by the report layer and counted as forbidden actions by the scorer.
\item \textbf{Uncertain outcomes.} A call whose result is unknown after a failure is never retried blindly. The
      attempt is retained and the run ends, because a repeated dispatch could duplicate a side effect.
\item \textbf{Host sandbox.} Real-model trials run in a restricted host environment (Section~\ref{sec:sandbox}).
\end{itemize}

\section{Host Sandbox for Real-Model Runs}
\label{sec:sandbox}
The model cannot run inside the container sandbox used on Linux, because MLX needs the host GPU. Real-model
trials therefore run the investigator from an installed wheel, in a fresh data root that holds only the visible
case, under a macOS \code{sandbox-exec} profile. The profile denies reads outside the visible, runtime and
virtual-environment paths and denies network access other than the loopback model endpoint. The environment
is stripped. Before the investigator starts, the launcher records probes (list and read of forbidden paths,
writes outside the root, a network connection, an import of evaluator code and an environment probe) and
refuses to start if any probe is not denied. In every real-model run reported here, all probes were denied
(Section~\ref{sec:reliab}). This is a host-level restriction and not a security proof, and the report
discloses that the operating-system profile is a deprecated mechanism.

\section{User Interface}
The browser interface consumes the HTTP API and does not compute compliance figures. It has screens to
connect to the API, to select a dossier, to start a run with a chosen variant, to follow an investigation as it
proceeds, to read the report and to record a reviewer decision, and a readiness screen that summarises the
state of the evidence and the checks. The investigation screen shows the sequence of turns with the plan,
the tool call and its arguments, the observation, and the assessment, so that a viewer can watch the agent plan,
act, observe and adapt.

\pending{Screenshots of the interface for Appendix~\ref{app:traces} have not yet been captured.}

\section{Repository Structure}
\begin{lstlisting}[language={},numbers=none,caption={Repository layout (abridged)},label=lst:repo]
src/recycleguard/         backend package
  intake/ extraction/ normalization/   evidence in
  storage/                             blob store, audit kernel
  indexing/ retrieval/ knowledge/      FTS5 index, retrieval, corpus
  domain/                              exact reconciler
  workflow/                            LangGraph runner, store, gate
  investigation.py local_model.py      tools and model gateway
  reporting/ review/ exports/          report, verifier, review, bundle
  api/                                 FastAPI application
benchmarks/               generator, oracle, scorer, trial coordinator
tools/local-model/        MLX server (separate, isolated dependencies)
ui/                       React and TypeScript interface
docs/                     design, contracts, evaluation protocol, checkpoints
tests/                    unit, integration and fixture dossiers
\end{lstlisting}
